\documentstyle[11pt]{article}
\setlength{\textwidth}{170mm}
\setlength{\textheight}{220mm}
\setlength{\oddsidemargin}{0mm}
\setlength{\evensidemargin}{10mm}
\setlength{\topmargin}{-10mm}
%\def\RR{\hbox{\raise.4ex\hbox{$\scriptstyle{|}$}\hskip-1.85pt R}}
\newcommand{\REF}{\nopagebreak\noindent{\large{\bf
      References}}\vspace{.2in}}

\newcommand{\Ref}[1]{\medskip\noindent\hangafter=1
      \hangindent=20pt\ignorespaces #1}
\pagestyle{empty}
\begin{document}


\noindent
University of Illinois \hfill Department of Economics \\
Spring 2002 \hfill Roger Koenker

%\vspace{4mm}
\begin{center}
\large{
{\bf Economics 476 \\
Term Paper Assignment}}
\end{center}

\vspace{5mm}
\begin{quotation}

``You haven't told me yet,'' said Lady Nuttal, ''What it is your fiance does for a living.''

``He's a statistician,'' replied Lamia, with an annoying sense of being on the defensive.

Lady Nuttal was obviously taken aback.  It had not occurred to her that statisticians 
entered into normal social relationships.  The species, she would have surmised, was perpetuated in some collateral manner, like mules.

``But Aunt Sara, it's a very interesting profession,'' said Lamia warmly.

``I don't doubt it,'' said her aunt, who obviously doubted it very much.  ``To 
express anything important in mere figures is so plainly impossible 
that there must be endless scope for well-paid advice on how to do it.  But 
don't you think that life with a statistician would be rather, shall we say, humdrum?''

Lamia was silent.  She felt reluctant to discuss the surprising depth of 
emotional possibility which she had discovered below Edward's numerical veneer.

``It's not the figures themselves,'' she said finally, ``It's what you do with them that matters.''

\hspace{45mm} K.A.C. Mandeville, {\em The Undoing of Lamia Gurdleneck}

\end{quotation}

\vspace{10mm}

A brief (5-15 pages) term paper will be expected from each student by 
Noon, Monday, May 4.  The paper may take one of three forms:

\begin{itemize}
\item a theoretical note on some aspect of econometrics we have touched upon 
in class, an elaboration of one of the problems for example.

\item an empirical note which explores the application of a technique we have 
considered on a ``real'' data set.

\item a critical survey of recent developments in a relevant niche of the 
econometrics literature, the emphasis here should be on the word 
{\em critical}.  You should try to offer some guidance on where the literature 
should go from here, not simply summarize in an undigested fashion.
\end{itemize}

While I don't encourage {\em purely} ``monte carlo'' papers, some monte-carlo 
{\em aspect} is encouraged.  To avoid potential misunderstandings about paper
topics, I would like all students to submit a 1-page ``paper proposal by 
March 9.  These proposals should contain at least three relevant references to 
literature and a brief sketch of what you propose to do in your paper.

I would strongly advise you to begin to learn \LaTeX and Splus as part of
the paper writing process. Some hopefully helpful hints on this are
available in my paper ``Reproducible Econometrics'' which I will distribute
in class.  I would also encourage you to buy William Thomson's "A
Guide for the Young Economist", published by MIT Press.  It contains much
good advice about writing papers, giving talks, and even writing referee
reports.  Thomson is a well respected theorist, but virtually all of his
advice is easily adapted to writing econometrics as well as economic theory.

\end{document}

